\section{Whole relative ideal restriction under genuine smooth submersions}

\begin{theorem}[The entire ideal of a simultaneous coordinate zero locus]
\label{hir:submersion-relative-coordinate-ideal}
\lean{Hironaka.HironakaSmoothSubmersionRelativeRestriction.vanishingIdeal_mem_iff_coordinate_family}
\uses{hir:submersion-relative-coordinate-families,bex:hadamard,bcalc:snc-simultaneous}
Let $A$ contain $x$ in a complex manifold of dimension $n$. Let $\Phi$ be a
genuine centered local holomorphic coordinate equivalence at $x$, and let
$S\subset\{0,\ldots,n-1\}$ be finite. Suppose that, on its actual chart domain,
$A$ is exactly the locus where all coordinates indexed by $S$ vanish. Then the
whole vanishing ideal at $x$ is generated by those actual coordinate functions.
This includes $S=\varnothing$, whose ideal is the zero ideal.
\leanok
\begin{proof}
Choose a genuine finite permutation taking $S$ to the last $|S|$ coordinates.
Its complex linear isometry and holomorphic inverse compose with $\Phi$ to
give another genuine centered chart. The same original set $A$ is now exactly
the coordinate plane of dimension $n-|S|$. Hadamard's lemma in this chart
expresses every holomorphic germ vanishing on $A$ as the sum of the last
$|S|$ coordinate functions with holomorphic coefficients; all earlier
coefficients are literally zero. Compose back with the inverse permutation.
Every surviving coordinate is one of the original coordinates indexed by
$S$, so this is a finite holomorphic span of the claimed actual functions.
Conversely each such coordinate is holomorphic and vanishes on $A$ on the
actual chart domain. Every finite holomorphic combination consequently
vanishes on $A$ near $x$. Empty $S$ leaves the empty sum, so the same argument
gives the literal zero germ without a nonemptiness assumption.
\end{proof}
\end{theorem}

\begin{theorem}[Actual finite presentations pull back in arbitrary dimensions]
\label{hir:submersion-relative-finite-presentation}
\lean{Hironaka.HironakaSmoothSubmersionRelativeRestriction.pullback_mem_iff_finite_family}
Let $f:Y^m\to X^n$ be holomorphic. Suppose the entire germ ideal $I$ at $f(y)$
is presented by a genuine finite family of holomorphic functions $g_i$.
Then $f^*I$ at $y$ is exactly the holomorphic span of the actual functions
$g_i\circ f$, with no equality of ambient dimensions required.
\leanok
\begin{proof}
By definition a member of $f^*I$ is a holomorphic combination of compositions
$a\circ f$, with $a\in I_{f(y)}$. Substitute the actual finite presentation
of each $a$; continuity pulls its whole germ identity back to a neighborhood
of $y$, and composition preserves the holomorphic coefficients. Transitivity
of finite holomorphic spans gives the claimed span of $g_i\circ f$.
Conversely every $g_i$ belongs to $I_{f(y)}$ by its own finite presentation,
so its composition belongs to the defining generating family of $f^*I$.
Closure under finite holomorphic combinations gives the reverse inclusion.
The empty presentation is handled by the same empty-span argument.
\end{proof}
\end{theorem}

\begin{theorem}[The whole original submanifold ideal pulls back]
\label{hir:submersion-relative-vanishing-ideal}
\lean{Hironaka.HironakaSmoothSubmersionRelativeRestriction.pullback_vanishingIdeal_submersion}
\uses{hir:submersion-relative-coordinate-ideal,hir:submersion-relative-finite-presentation,hir:submersion-relative-coordinate-families,hir:submersion-relative-regular-geometry}
For a genuine holomorphic submersion $f:Y^m\to X^n$ and an original closed
complex submanifold $N\subset X$, the entire germ ideals satisfy
$f^*\mathcal I_N=\mathcal I_{f^{-1}N}$.
\leanok
\begin{proof}
Fix $y$. If $f(y)\in N$, choose the actual original flattening chart supplied
by regularity of $N$. Its last $n-j$ coordinates define the whole original
set $N$ locally. The genuine submersion chart construction internally lifts
the entire coordinate family to one centered chart at $y$, with an actual
embedding of all original coordinate indices. Thus $f^{-1}N$ is exactly the
common zero locus of the lifted last coordinates on this actual chart.
The first theorem gives the entire ideal on both manifolds. The second
theorem identifies the pullback with the span of the composed original
coordinates. The actual chart dictionary identifies these functions as germs
with the lifted coordinates, so the spans, and therefore the whole ideals,
are equal. If $f(y)\notin N$, closedness gives an actual neighborhood disjoint
from $N$ and its inverse image disjoint from $f^{-1}N$. Both ideals are the
unit ideal there; the genuine generator $1$ pulls back to $1$. Pointwise
equality for every holomorphic function and every $y$ gives equality of the
entire germ-ideal records. No caller flattening or ideal-equality certificate
is supplied.
\end{proof}
\end{theorem}

\begin{theorem}[Whole sums pull back across ambient dimensions]
\label{hir:submersion-relative-ideal-sums}
\lean{Hironaka.HironakaSmoothSubmersionRelativeRestriction.pullback_sumIdeal_dimensions}
For a holomorphic map $f:Y^m\to X^n$ and original coherent germ ideals $I,J$,
$f^*(I+J)=f^*I+f^*J$ as entire germ-ideal records.
\leanok
\begin{proof}
A member of the left side is a holomorphic combination of compositions of
members of $I+J$. Expand each such member as a finite holomorphic combination
of members of $I$ or $J$, pull the genuine germ identity back by continuity,
and compose its holomorphic coefficients. Transitivity gives membership in
the sum on the right. Conversely a composed member of either original ideal
is a composed member of $I+J$, because coherence makes the original members
holomorphic and the sum contains both ideals. Every generator of the sum on
the right therefore belongs to the pullback on the left, and transitivity of
the defining finite holomorphic spans proves the reverse inclusion. These
constructions use the actual dimensions $m$ and $n$ throughout.
\end{proof}
\end{theorem}

\begin{theorem}[The complete original restricted ideal dictionary]
\label{hir:submersion-relative-restricted-ideal}
\lean{Hironaka.HironakaSmoothSubmersionRelativeRestriction.pullback_restrictOn_submersion}
\uses{hir:submersion-relative-vanishing-ideal,hir:submersion-relative-ideal-sums,ana:restricted-calculus}
For the same genuine submersion, original closed submanifold $N$ and coherent
original ideal $I$, the ambient representative of its restriction satisfies
$f^*\operatorname{Res}_N I=\operatorname{Res}_{f^{-1}N}(f^*I)$ as an entire
germ-ideal record.
\leanok
\begin{proof}
The actual restricted-ideal construction is $\operatorname{Res}_N I=
I+\mathcal I_N$: subtract the finite restricted-span identity to obtain a
holomorphic remainder vanishing on $N$, and conversely add such a remainder
to any original finite span. Apply the fourth theorem to this genuine sum,
and the third theorem to the whole original vanishing ideal. The pullback of
the coherent original ideal is coherent, so the same actual restricted-ideal
identity on $Y$ identifies the resulting sum with the right side. Every
operation concerns the original whole records and the actual original set.
\end{proof}
\end{theorem}

\begin{theorem}[Coherence and the entire restricted marked record]
\label{hir:submersion-relative-restricted-marked-record}
\lean{Hironaka.HironakaSmoothSubmersionRelativeRestriction.isCoherent_restrictOn_preimage_submersion,Hironaka.HironakaSmoothSubmersionRelativeRestriction.markedPullback_restricted_record}
\uses{hir:submersion-relative-restricted-ideal,hir:submersion-relative-regular-geometry,ana:restricted-calculus}
The whole pulled restricted ideal above is coherent. Replacing the ideal in
an original marked record by its actual restriction commutes with the entire
marked pullback, retaining the original marking and every original ordered
boundary entry.
\leanok
\begin{proof}
The original submanifold ideal is coherent by its actual finite flattening
coordinates, including the genuine unit presentation off the submanifold.
Its sum with the coherent original ideal is coherent. Holomorphic pullback
preserves coherence, and the fifth theorem identifies this actual coherent
ideal with the claimed pulled restricted ideal. For the marked record, the
underlying set on both sides is literally $f^{-1}N$, the ordered boundary is
the same original list mapped entry by entry by inverse image, and the
marking is unchanged. The fifth theorem equates the whole ideal field.
Equality of all fields consequently gives equality of the entire marked
records, including empty and repeated boundary entries. No higher invariant
equality or general blowup base change is asserted here; those parts of G2
remain open.
\end{proof}
\end{theorem}
